\section*{अध्याय \ref{ch:b2:reproductive-hormones} --- जनन का हॉर्मोनी नियंत्रण}

\begin{solution}{exo:b2:reproductive-hormones:1}
अधश्चेतक (GnRH, स्पंदों में) $\to$ अग्र पीयूष (LH, FSH) $\to$
\omterm{def:b2:mammal-reproduction:testis}{वृषण}: LH \omterm{def:b2:mammal-reproduction:testis}{लाइडिग कोशिकाओं} पर (टेस्टोस्टेरोन), FSH \omterm{def:b2:mammal-reproduction:testis}{सर्टोली कोशिकाओं} पर
(\omterm{def:b2:mammal-reproduction:testis}{शुक्राणुजनन}, इनहिबिन)। प्रतिपुष्टि: टेस्टोस्टेरोन GnRH और
LH को संदमित करता है; इनहिबिन FSH को।
\end{solution}

\begin{solution}{exo:b2:reproductive-hormones:2}
अंडाशय: \omterm{def:b2:reproductive-hormones:cycle}{पुटकीय प्रावस्था}, दिन 1--14, बढ़ते \omterm{def:b2:mammal-reproduction:ovary}{पुटक} का एस्ट्राडायोल;
\omterm{def:b2:mammal-reproduction:ovary}{अंडोत्सर्ग}, दिन 14, \omterm{def:b2:reproductive-hormones:cycle}{LH उछाल}; \omterm{def:b2:reproductive-hormones:cycle}{पीत प्रावस्था}, दिन 15--28,
\omterm{def:b2:mammal-reproduction:ovary}{पीत पिंड} का प्रोजेस्टेरोन। गर्भाशय: रजःस्राव, दिन 1--5
(स्टेरॉइड-वापसी); प्रसारी प्रावस्था, दिन 6--14 (एस्ट्राडायोल);
स्रावी प्रावस्था, दिन 15--28 (प्रोजेस्टेरोन)।
\end{solution}

\begin{solution}{exo:b2:reproductive-hormones:3}
दैनिक एस्ट्रोजन और प्रोजेस्टिन ऋण प्रतिपुष्टि से LH तथा FSH नीची रखते
हैं, इसलिए कोई \omterm{def:b2:mammal-reproduction:ovary}{पुटक} भर्ती नहीं होता, कोई एस्ट्राडायोल-चढ़ाई नहीं होती और
वह उछाल कभी नहीं छूटती; और वह प्रोजेस्टिन श्लेष्म भी गाढ़ा करता है।
गोली-रहित सप्ताह में वे स्टेरॉइड गिरते हैं और उन्हीं के नीचे बना
\omterm{def:b2:reproductive-hormones:cycle}{गर्भाशय अंतःस्तर} झड़ जाता है --- यानी वापसी-रक्तस्राव, न कि किसी
\omterm{def:b2:mammal-reproduction:ovary}{अंडोत्सर्ग} के बाद का रजःस्राव।
\end{solution}

\begin{solution}{exo:b2:reproductive-hormones:4}
LH और FSH कई गुना चढ़ जाती हैं (न टेस्टोस्टेरोन की प्रतिपुष्टि, न
इनहिबिन की); पेशी, कामेच्छा और सहायक ग्रंथियाँ क्षीण हो जाती हैं।
टेस्टोस्टेरोन उन लक्षणों को लौटा देता है और LH फिर नीचे ले आता है, पर
उर्वरता नहीं: \omterm{def:b2:mammal-reproduction:testis}{वृषण} जा चुके हैं, और \omterm{def:b2:mammal-reproduction:testis}{वृषण} वाले किसी पुरुष में भी बाहर से
दिया गया टेस्टोस्टेरोन LH दबा देता है और इसलिए उस वृषण-भीतरी
टेस्टोस्टेरोन को भी जो \omterm{def:b2:mammal-reproduction:testis}{शुक्राणुजनन} को चाहिए।
\end{solution}

\begin{solution}{exo:b2:reproductive-hormones:5}
प्रतिदिन $24/1.5 = 16$ स्पंदें; और $24/4 = 6$। जब \omterm{def:b2:mammal-reproduction:ovary}{पीत पिंड} मरता है
तब स्पंदें अब भी धीमी हैं, जो FSH के पक्ष में है: FSH पहले चढ़ती है ---
और वही सही हॉर्मोन है, क्योंकि \omterm{def:b2:mammal-reproduction:ovary}{पुटकों} का अगला दल FSH ही
भर्ती करती है।
\end{solution}

\begin{solution}{exo:b2:reproductive-hormones:6}
\qty{200}{pg/mL} $= \qty{200}{ng/L}$: $200\times 10^{-9}/272 =
\qty{7.4e-10}{mol/L} = \qty{740}{pmol/L}$। \qty{15}{ng/mL} $=
\qty{15}{\micro g/L}$: $15\times 10^{-6}/314 = \qty{48}{nmol/L}$। और
एक माइक्रोलिटर में: $7.4\times 10^{-16}$ mol $= 4.4\times 10^{8}$
अणु।
\end{solution}

\begin{solution}{exo:b2:reproductive-hormones:7}
\omterm{def:b2:mammal-reproduction:ovary}{अंडोत्सर्ग} दिन $35 - 14 = 21$ पर; उर्वर दिन 18 से दिन 22 तक। अनियमित
चक्रों में \omterm{def:b2:reproductive-hormones:cycle}{पुटकीय प्रावस्था}, और इसलिए \omterm{def:b2:mammal-reproduction:ovary}{अंडोत्सर्ग} का दिन,
पिछले चक्रों से भाँपा नहीं जा सकता, और वह खिड़की छूट जाती है या ग़लत
जगह पड़ जाती है।
\end{solution}

\begin{solution}{exo:b2:reproductive-hormones:8}
दिन 21 से 26 तक 2.5 दुगुनेपन हैं: $2^{2.5} = 5.7$ IU/L --- यानी पर्याप्त।
दिन 25 पर \omterm{def:b2:mammal-reproduction:implantation}{रोपण} हो तो दिन 26 पर $1.4$ IU/L मिलती है; और \omterm{def:b2:mammal-reproduction:ovary}{पीत पिंड},
जो पहले ही क्षीण हो रहा है, बचता नहीं, प्रोजेस्टेरोन गिरता है और वह
भ्रूण रजःस्राव के साथ खो जाता है। देर से \omterm{def:b2:mammal-reproduction:implantation}{रोपण} आरंभिक हानि का
एक आम कारण है।
\end{solution}

\begin{solution}{exo:b2:reproductive-hormones:9}
पीयूष GnRH के ढंग को उत्तर देता है, उसकी मात्रा को नहीं:
प्रति घंटा स्पंदें LH तथा FSH बनाए रखती हैं, जबकि वही दैनिक मात्रा
लगातार डाली जाए तो उन्हें मिटा देती है। इसने किसी सरल मात्रा--अनुक्रिया
को ख़ारिज कर दिया। महीने भर के दीर्घक्रिय प्रचालक पर: कुछ दिनों की भड़क
के बाद LH गिरकर लगभग शून्य हो जाती है (\omterm{prop:b2:reproductive-hormones:pulses}{ग्राही असंवेदन}), एस्ट्राडायोल
रजोनिवृत्ति के स्तरों पर, \omterm{def:b2:reproductive-hormones:cycle}{गर्भाशय अंतःस्तर} क्षीण हो जाता है और कोई चक्र
नहीं होता --- यानी एक प्रतिवर्ती चिकित्सकीय रजोनिवृत्ति।
\end{solution}

\begin{solution}{exo:b2:reproductive-hormones:10}
बिना हॉर्मोन: $R^* = 100/0.5 = 200$। सतत: $100/6.5 = 15$। कोई
दो-मिनट की स्पंद $6\times 200\times(2/60) = 40$ ग्राही हटाती है
(\qty{20}{\%}); और अगले घंटे में वह घाटा
$e^{-0.5}$ गुणक से सिकुड़ता है, यानी उसका \qty{39}{\%} उबर जाता है। वह कोष
वहाँ ठहरता है जहाँ प्रति स्पंद हानि प्रति घंटा उबार के बराबर हो:
यानी हर स्पंद से पहले लगभग 150 ग्राही, अधिकतम का तीन चौथाई --- यानी
सतत स्तर का दस गुना।
\end{solution}

\begin{solution}{exo:b2:reproductive-hormones:11}
रजोनिवृत्ति: कोई \omterm{def:b2:mammal-reproduction:ovary}{पुटक} नहीं, इसलिए न एस्ट्राडायोल न इनहिबिन, इसलिए कोई
प्रतिपुष्टि नहीं: LH और FSH ऊँची चलती हैं। प्रोलैक्टिनोमा: प्रोलैक्टिन
GnRH स्पंद-जनित्र धीमा कर देती है, इसलिए LH कम है, \omterm{def:b2:mammal-reproduction:ovary}{पुटक} चलाए नहीं जाते,
एस्ट्राडायोल कम है और कोई चक्र नहीं। कणिकामय-कोशिका अर्बुद:
लगातार ऊँचा एस्ट्राडायोल ऋण प्रतिपुष्टि से LH दबा देता है (और
बिना किसी पकते \omterm{def:b2:mammal-reproduction:ovary}{पुटक} के \omterm{def:b2:mammal-reproduction:ovary}{अंडोत्सर्ग} को कुछ है ही नहीं), इसलिए
चक्र ऊँचे एस्ट्राडायोल के साथ रुक जाता है --- और \omterm{def:b2:reproductive-hormones:cycle}{गर्भाशय अंतःस्तर} बढ़ता चला जाता है।
\end{solution}

\begin{solution}{exo:b2:reproductive-hormones:12}
कम या थोड़ी देर का एस्ट्राडायोल GnRH और LH घटाता है; जबकि ऊँचा, टिकाऊ
एस्ट्राडायोल (\qty{200}{pg/mL} से ऊपर 36 घंटे से अधिक)
किसपेप्टिन--GnRH तंत्र की और पीयूष की अनुक्रिया बदल देता है,
और पीयूष इस बीच LH भर चुका होता है, इसलिए वही हॉर्मोन अब
एक उछाल भड़का देता है। \omterm{def:b2:mammal-reproduction:ovary}{अंडोत्सर्ग} को सर्वथा-अथवा-कुछ नहीं और समय-सधा
होना ही चाहिए: वह अंडाणु एक ही बार, एक ही दिन, तब छूटना चाहिए जब कोई
\omterm{def:b2:mammal-reproduction:ovary}{पुटक} पका हो। कोई डायल
(एस्ट्राडायोल के समानुपाती LH) क्रमिक, आंशिक
पीतीकरण देता, \omterm{def:b2:mammal-reproduction:ovary}{पुटक} आधे-पके या कई एक साथ अंडोत्सर्जित होते,
और कोई निश्चित उर्वर दिन ही न होता।
\end{solution}

\begin{solution}{pb:b2:reproductive-hormones:1}
\textbf{1.} पुटकीय, दिन 1--14: एस्ट्राडायोल चढ़ता हुआ, प्रोजेस्टेरोन
कम। पीत, दिन 15--28: प्रोजेस्टेरोन ऊँचा।
\textbf{2.} वह उछाल दिन 14 पर शिखर पर थी (आरंभ दिन 13 को);
\omterm{def:b2:mammal-reproduction:ovary}{अंडोत्सर्ग} उसके आरंभ के लगभग 36 घंटे बाद, दिन 15 पर।
\textbf{3.} दिन 15 से 28: यानी 14 दिन, वही \omterm{def:b2:mammal-reproduction:ovary}{पीत पिंड} की
आयु।
\textbf{4.} प्रोजेस्टेरोन तापनियमन का नियत बिंदु
\qty{0.3}{\celsius} से \qty{0.4}{\celsius} चढ़ा देता है; और वह चढ़ाई
\omterm{def:b2:mammal-reproduction:ovary}{अंडोत्सर्ग} के एक-दो दिन बाद दिखती है, इसलिए वह पुष्टि करती है कि
\omterm{def:b2:mammal-reproduction:ovary}{अंडोत्सर्ग} हो चुका है पर उसकी घोषणा नहीं कर सकती।
\textbf{5.} दिन 11 से 14 (210, 260, 310, 220): यानी तीन से चार दिन,
और वह स्विच जिन 36 घंटों की माँग करता है उससे अधिक।
\textbf{6.} दिन 12 से 16।
\textbf{7.} $310\times 10^{-9}/272 = \qty{1.14e-9}{mol/L} =
\qty{1140}{pmol/L}$; $16\times 10^{-6}/314 = \qty{51}{nmol/L}$।
\textbf{8.} $65/8 = 8$। एक घंटे के अर्ध-जीवन के साथ, स्राव रुकते ही LH
हर घंटे आधी हो जाती है, इसलिए दिन भर बाद वह फिर आधाररेखा के पास है।
\textbf{9.} \omterm{def:b2:mammal-reproduction:ovary}{अंडोत्सर्ग} लगभग दिन 22 पर, और चक्र लगभग 36 दिन का।
\textbf{10.} \omterm{def:b2:mammal-reproduction:implantation}{रोपण} दिन 22 पर; और दिन 27 तक (पाँच दिन, 2.5
दुगुनेपन) लगभग \qty{5.7}{IU/L}।
\textbf{11.} दिन 27 पर प्रोजेस्टेरोन का गिरकर \qty{2}{ng/mL} होना
बताता है कि \omterm{def:b2:mammal-reproduction:ovary}{पीत पिंड} समय पर क्षीण हो रहा है: उसे किसी ने बचाया नहीं।
\textbf{12.} प्रोजेस्टेरोन लगभग \qty{20}{ng/mL} और चढ़ता हुआ, तथा
एस्ट्राडायोल लगभग \qty{200}{pg/mL}।
\textbf{13.} चक्र के अंत में एस्ट्राडायोल, प्रोजेस्टेरोन और
इनहिबिन साथ-साथ गिरते हैं और FSH अपनी प्रतिपुष्टि से बच निकलती है; फिर
जैसे-जैसे वह दल बढ़ता है, एस्ट्राडायोल तथा इनहिबिन चढ़ते हैं और FSH
गिरती है, और गिरती FSH पर केवल वही \omterm{def:b2:mammal-reproduction:ovary}{पुटक} बढ़ता रहता है जिसके FSH ग्राही
(यानी कणिकामय कोशिकाएँ) सबसे अधिक हैं --- बाक़ी मर जाते हैं।
\textbf{14.} बिना हॉर्मोन $R^* = 200$; और सतत हॉर्मोन के अंतर्गत
$100/6.5 = 15$।
\textbf{15.} $6\times 200\times 2/60 = 40$ ग्राही, यानी \qty{20}{\%};
काल-नियतांक $1/d = \qty{2}{h}$, इसलिए घंटे भर में उस घाटे का \qty{39}{\%}
उबर जाता है; और वह कोष लगभग 150 ग्राही पर ठहरता है, यानी अधिकतम का तीन
चौथाई।
\textbf{16.} स्पंदें पीयूष को लगभग 150 ग्राही पर रखती हैं और वह
हर स्पंद का उत्तर LH से देता है; जबकि कोई सतत निवेश उस कोष को
15 तक ले जाता है और वही कोशिकाएँ चुप हो जाती हैं, मात्रा चाहे कुछ भी हो;
और स्पंदें उस कोष तथा उस अनुक्रिया को फिर बहाल कर देती हैं।
\textbf{17.} पहले दिन: LH और टेस्टोस्टेरोन चढ़ते हैं (वह प्रचालक उन
ग्राहियों को उद्दीपित करता है जो अब भी हैं --- यानी भड़क)। एक महीने बाद:
ग्राही असंवेदित, LH लगभग शून्य, टेस्टोस्टेरोन बधिया
स्तरों पर, और यही चिकित्सकीय लक्ष्य है।
\textbf{18.} FSH चढ़ती है (कोई इनहिबिन नहीं), LH सामान्य रहती है;
हर चक्र में अधिक \omterm{def:b2:mammal-reproduction:ovary}{पुटक} चुनाव झेल जाते हैं, और कई \omterm{def:b2:mammal-reproduction:ovary}{अंडोत्सर्ग} तथा
द्वियुग्मज जुड़वाँ होते हैं।
\textbf{19.} किसी प्राकृतिक चक्र में FSH की गिरावट प्रबल \omterm{def:b2:mammal-reproduction:ovary}{पुटक} के
सिवा सबको मार देती है; जबकि FSH की अचर आपूर्ति पूरे दल को
परिपक्वता तक बढ़ने देती है।
\textbf{20.} दस \omterm{def:b2:mammal-reproduction:ovary}{पुटक} जल्दी ही \qty{200}{pg/mL} से कहीं अधिक
एस्ट्राडायोल बनाते हैं, जो कोई असमय \omterm{def:b2:reproductive-hormones:cycle}{LH उछाल} और संग्रह से
पहले ही \omterm{def:b2:mammal-reproduction:ovary}{अंडोत्सर्ग} करा देता; और वह प्रतिरोधी GnRH रोक देता है, इसलिए कोई
उछाल हो ही नहीं सकती।
\textbf{21.} सारणी में वह उछाल दिन 14 पर शिखर पर थी और \omterm{def:b2:mammal-reproduction:ovary}{अंडोत्सर्ग}
उसके आरंभ के लगभग 36 घंटे बाद हुआ; hCG उस उछाल की जगह लेती है,
और 35 घंटे पर संग्रह अंडाणुकों को परिपक्व पर अब भी
\omterm{def:b2:mammal-reproduction:ovary}{पुटक} के भीतर पकड़ लेता है।
\textbf{22.} $10\times 0.7 = 7$ भ्रूण, $7\times 0.4 = 2.8$
कोरकपुटियाँ; $1 - 0.65^{2} = 0.58$।
\textbf{23.} अचर एस्ट्रोजन और प्रोजेस्टिन FSH तथा LH नीची रखते हैं: कोई
\omterm{def:b2:mammal-reproduction:ovary}{पुटक} नहीं बढ़ता, एस्ट्राडायोल कभी उस देहली तक नहीं चढ़ता, कोई उछाल नहीं,
कोई \omterm{def:b2:mammal-reproduction:ovary}{पीत पिंड} नहीं, कोई प्रोजेस्टेरोन-चढ़ाई नहीं।
\textbf{24.} गौण लक्षणों और कामेच्छा के लिए टेस्टोस्टेरोन देना ही पड़ेगा;
पर वह उस ऊँची वृषण-भीतरी सांद्रता को नहीं लौटाता
जो \omterm{def:b2:mammal-reproduction:testis}{लाइडिग कोशिकाएँ} देती थीं, इसलिए \omterm{def:b2:mammal-reproduction:testis}{शुक्राणुजनन} दबा ही
रहता है --- और मर्म यही है।
\textbf{25.} \omterm{def:b2:mammal-reproduction:ovary}{अंडोत्सर्ग} दिन 15 पर; \omterm{def:b2:reproductive-hormones:cycle}{पीत प्रावस्था} 14 दिन; उर्वर
खिड़की दिन 12--16; और वह उछाल एस्ट्राडायोल के
\qty{200}{pg/mL} से ऊपर 36 घंटे से अधिक रहने से चली, जो दिन 11--14 पर पूरा हुआ।
\end{solution}
